% ============================================================
% CONCLUSION GÉNÉRALE ET PERSPECTIVES
% ============================================================

\chapter*{Conclusion Générale et Perspectives}
\addcontentsline{toc}{chapter}{Conclusion Générale et Perspectives}
\markboth{Conclusion Générale et Perspectives}{Conclusion Générale et Perspectives}

\vspace{0.5cm}

\section*{Bilan du projet}

Le présent travail de fin d'études, réalisé au sein de \textbf{TIM Tunisie}, avait pour objectif de concevoir et de développer une plateforme SaaS de gouvernance et de monitoring des modèles de langage de grande taille (LLM). Au terme de trois sprints de développement conduits selon la méthodologie Scrum, nous avons livré \textbf{LLM Dashboard}, une solution intégrée et opérationnelle répondant de manière exhaustive à la problématique identifiée.

La plateforme développée se distingue par la richesse et la complémentarité de ses \textbf{onze modules fonctionnels}~:

\begin{itemize}[noitemsep]
    \item Le \textbf{Gateway LLM} assure le proxying transparent des appels vers cinq fournisseurs (OpenAI, Anthropic, Google, Groq, Cohere), avec journalisation complète dans TimescaleDB et évaluation des règles de gouvernance en temps réel~;
    \item Le \textbf{moteur de gouvernance} applique cinq types de règles configurables (quota, budget, restriction de modèle, limite de tokens, restriction temporelle) avec des compteurs atomiques Redis~;
    \item Le \textbf{pipeline de détection d'anomalies} combine trois algorithmes de machine learning complémentaires --- STL decomposition~\cite{cleveland1990}, Isolation Forest~\cite{liu2008} et CUSUM~\cite{page1954} --- pour une couverture maximale des types d'anomalies~;
    \item Le \textbf{module Explainer} génère automatiquement des explications intelligibles et des recommandations actionnables à l'aide de LLM, avec stockage des embeddings vectoriels pour la recherche sémantique~;
    \item Le \textbf{moteur de prévision budgétaire} projette les coûts futurs selon trois scénarios (optimiste, moyen, pessimiste) avec identification des risques de dépassement~;
    \item Le \textbf{système de facturation contractuelle} supporte trois types de contrats (pay-as-you-go, forfait, hybride) avec génération automatique de factures détaillées~;
    \item Le \textbf{tableau de bord interactif}, développé avec Next.js~16 et React~19, offre une visualisation en temps réel via WebSocket de toutes les métriques de la plateforme.
\end{itemize}

\vspace{0.3cm}

Sur le plan technique, la plateforme repose sur une \textbf{architecture trois-tiers} robuste~: un backend FastAPI asynchrone, une base de données PostgreSQL enrichie des extensions TimescaleDB et pgvector, un cache multi-niveaux Redis, un pipeline de tâches asynchrones Celery, et un déploiement conteneurisé via Docker Compose avec six services orchestrés. La qualité du code est validée par \textbf{180 tests} (124 unitaires et 56 d'intégration) et un pipeline CI/CD automatisé via GitHub Actions.

\vspace{0.3cm}

Du point de vue académique, ce projet a permis de mobiliser et d'approfondir des connaissances dans des domaines variés~: l'architecture logicielle (Clean Architecture, principes SOLID)~\cite{martin2018}, l'apprentissage automatique non supervisé pour la détection d'anomalies~\cite{chandola2009}, les architectures multi-tenant SaaS~\cite{bezemer2010}, les API REST et les protocoles de sécurité (JWT, PBKDF2, TLS)~\cite{rfc7519,owasp2021}, et les systèmes de traitement de séries temporelles~\cite{hyndman2018}.


\section*{Limites identifiées}

Malgré les résultats obtenus, plusieurs limites ont été identifiées~:

\begin{enumerate}[noitemsep]
    \item \textbf{Tests de charge}~: les tests de performance sous charge élevée (milliers de requêtes simultanées) n'ont pas été réalisés en conditions réelles de production~;
    \item \textbf{Seuils de détection}~: les seuils des algorithmes de détection d'anomalies sont actuellement statiques et nécessiteraient un calibrage adaptatif par tenant~;
    \item \textbf{Génération PDF}~: le format des factures PDF pourrait être enrichi avec un template graphiquement plus élaboré~;
    \item \textbf{Alertes multi-canal}~: le système de notification est fonctionnel mais limité au webhook~; l'intégration avec les canaux email et Slack reste à compléter~;
    \item \textbf{Modèles de prévision}~: les algorithmes de prévision budgétaire utilisent des méthodes linéaires et pourraient bénéficier de modèles plus sophistiqués (ARIMA, Prophet).
\end{enumerate}


\section*{Perspectives d'évolution}

Les perspectives d'évolution de la plateforme LLM Dashboard sont multiples et s'inscrivent dans une logique d'amélioration continue~:

\begin{enumerate}[noitemsep]
    \item \textbf{Auto-calibrage des seuils ML}~: implémenter un mécanisme d'apprentissage en ligne (\textit{online learning}) pour ajuster automatiquement les seuils de détection d'anomalies en fonction du comportement de chaque tenant~;
    \item \textbf{Prévision avancée}~: intégrer des modèles de prévision de séries temporelles plus sophistiqués tels que ARIMA~\cite{box2015}, Prophet ou des réseaux de neurones récurrents (LSTM)~;
    \item \textbf{Analyse de sentiment des prompts}~: exploiter les modèles NLP pour analyser le contenu des prompts et détecter les usages potentiellement problématiques~;
    \item \textbf{Marketplace de modèles}~: développer une place de marché interne permettant aux tenants de découvrir et de comparer les modèles LLM disponibles~;
    \item \textbf{Déploiement Kubernetes}~: migrer de Docker Compose vers Kubernetes pour une scalabilité horizontale automatique et une haute disponibilité~;
    \item \textbf{Conformité réglementaire}~: intégrer des contrôles de conformité RGPD et AI Act pour garantir une utilisation responsable des LLM~;
    \item \textbf{Extension multi-cloud}~: supporter le déploiement sur les trois principaux clouds (AWS, GCP, Azure) avec abstraction de l'infrastructure.
\end{enumerate}

\vspace{0.5cm}

En définitive, ce projet de fin d'études nous a permis de concevoir et de livrer une plateforme complète et opérationnelle qui répond aux besoins réels de TIM Tunisie en matière de gouvernance des LLM. L'expérience acquise, tant sur le plan technique que méthodologique, constitue un socle solide pour les évolutions futures de la solution et pour notre parcours professionnel dans le domaine de l'ingénierie logicielle et de l'intelligence artificielle.
